\chapter{The LeanFM Project Bootstrap Contract}
\label{ch:bootstrap-contract}

This chapter defines the minimum machine-readable context that makes a LeanFM
project self-hosting. It belongs in the machine reference because an agent must
apply it before it can turn the rest of the book into checked artifacts. The
bootstrap defines the language and its tools; it must not contain a hidden copy
of the product requirements or implementation.

The intended handoff is deliberately simple: send another person the URL of
this PDF, ask Claude, Gemini, Codex, or another capable construction agent to
read it, and ask that agent to make the software. The current generic bootstrap
distribution is published at
\url{https://github.com/rfielding/leanFM/tree/main/projects/blank-leanfm-webserver}.
An agent using only the PDF may follow that URL for the language files and
commands named below. It must not copy the existing LeanFM application source
or prior ChatUI specifications as an answer key.

\section{Required project inputs}

A copyable project root contains the following stable inputs:

\begin{description}[style=nextline]
  \item[\texttt{book/main.pdf}] Central product-design argument and maintained
        system-prompt context.
  \item[\texttt{book/main-dark.pdf}] Accessible rendering of the same source.
  \item[\texttt{LeanFM/}] Versioned generic authoring types and deterministic
        validators. Product actors, messages, tasks and routes do not belong in
        this directory.
  \item[\texttt{bootstrap/}] Programs that print the language references and
        interrogation checklist, validate the conventional specification, and
        derive review views from its typed value.
  \item[\texttt{lean-toolchain} and \texttt{lakefile.toml}] The exact compiler
        selection and module roots needed to run the bootstrap locally.
  \item[\texttt{project.json}] Project identity, stable artifact paths,
        acceptance state and the classification of design versus language
        inputs.
\end{description}

All of these files are below the exclusive project root. A project may use an
external immutable dependency cache, but no project-owned mutable state or
acceptance decision may live there.

The project root is also the complete input to a fresh agent conversation. The
agent discovers the stable paths from the manifest and repository instructions;
the stakeholder need not paste a transcript or restate a bootstrap procedure.
If the root has no durable problem statement, LeanFM creates the minimal
editable book shell and begins with the visible question ``What problem are you
trying to solve?'' The answer starts requirements discovery. Named ambiguity
questions, book revisions, validation evidence and the two acceptance decisions
then accumulate beneath that same root.

\section{Conventional generated interface}

The requirements candidate lives at
\texttt{Specification/Requirements.lean}, imports
\texttt{LeanFM.Artifacts}, and defines:

\begin{lstlisting}[language=Lean]
namespace LeanFMProject.Requirements
def spec : LeanFM.RequirementSpec := ...
def all : List LeanFM.GeneratedRequirement := [.requirement spec]
end LeanFMProject.Requirements
\end{lstlisting}

The byte vocabulary lives beside it at
\texttt{Specification/Requirements.proto}. After requirements acceptance, the
separate implementation mapping lives at
\texttt{Specification/Implementation.lean}. Capitalization is significant
because these are Lean module paths, not informal filenames.

\section{Required commands and order}

The project supplies deterministic equivalents of these commands:

\begin{lstlisting}[language=bash]
make reference       # print the complete requirements-language reference
make interrogation   # print questions that expose missing observations
make validate        # elaborate and run LeanFM plus protobuf validation
make render          # derive interaction and residual-state review views
make check           # validate, then render
\end{lstlisting}

An agent runs \texttt{reference} and \texttt{interrogation} before authoring.
It must not infer constructors from prose or invent a smaller Lean lookalike.
The validator consumes the exported \texttt{RequirementSpec}; successful Lean
elaboration by itself is not LeanFM validation. The renderer also consumes that
same value. A separately written diagram is not evidence that it matches the
requirements.

\section{Acceptance gates}

The first gate consists of the requirements value, protobuf schema, derived
views, ambiguity register and validation evidence. A model, successful build or
attractive diagram cannot accept it. After explicit stakeholder acceptance, the
second gate adds the traceable implementation specification. Application source
and tests may be written only after both gates.

The bootstrap therefore solves a circularity without weakening the gates. The
book explains what is being built; the generic distribution makes the LeanFM
language executable; accepted project artifacts preserve the decisions. None
is a substitute for the other two.

\section{Construction is not a runtime dependency}

Codex is not part of the delivered server. Claude, Gemini, Codex, or a human may
act as the construction agent that turns this PDF and the generic bootstrap
into accepted specifications and source. After construction, the user starts
the LeanFM server, visits its root URL in an ordinary browser, creates or logs
into an account, supplies a model-provider key, and conducts the requirements
argument through LeanFM's own ChatGPT-like interface. The running server talks
to the selected model provider through its documented message protocol; it does
not require the construction agent, its conversation, or its local tooling.

Consequently, a reconstruction is not successful merely because it can be
driven from a coding-agent terminal. It must build a standalone executable,
serve the complete browser workflow, start without a platform API key, isolate
concurrent users and their credentials, and pass its conformance tests without
Codex being present.

\section{Local and server-hosted references}

The local commands are authoritative when no server is running. A running
LeanFM server additionally exposes the requirements and implementation language
references, the interrogation checklist, and read-only MCP context without
requiring a user login. These public endpoints contain language documentation,
not project or account data. Project artifacts, conversations, credentials and
acceptance mutations remain authenticated and owner-scoped.

